% Provenance (never printed): the labour strand of ../economics_series.md (grade 8: the labour market; employed,
% unemployed, outside; the two rates and the survey; why wages differ; wages and productivity over long periods).
% Brief: sources/book1/BRIEFS.md, g8-05. Chart data: figdata/grade-8/05-jobs-wages-unemployment (JST R6 unemp, wage,
% cpi, rgdpmad; US 2019 from WDI SL.UEM.TOTL.NE.ZS because JST prints 2.0; Brenland accounts and households).
% Ledger: g8:bls-unemployed-search, g8:bls-labour-force-age, g8:ilo-future-starters.
\chapter{Jobs, Wages and Unemployment}\label{ch:g8:jobs-wages-unemployment}

\begin{openingchart}{Unemployment rate in five countries, 1970--2020, in percent of the labour
force.}{Jord\`a--Schularick--Taylor Macrohistory Database, \omterm{def:g8:jobs-wages-unemployment:unemployment-rate}{unemployment rate}; United States 2019
from World Bank, World Development Indicators, SL.UEM.TOTL.NE.ZS}
\begin{tikzpicture}
\begin{axis}[omtime, xmin=1970, xmax=2020, ymin=0, ymax=28, xlabel={Year},
  ylabel={Unemployment rate (percent)},
  legend style={at={(0.5,-0.2)}, anchor=north, legend columns=5}]
\addplot table[x=year, y=usa] {figdata/grade-8/05-jobs-wages-unemployment/out/unemployment.csv};
\addplot table[x=year, y=gbr] {figdata/grade-8/05-jobs-wages-unemployment/out/unemployment.csv};
\addplot table[x=year, y=deu] {figdata/grade-8/05-jobs-wages-unemployment/out/unemployment.csv};
\addplot table[x=year, y=esp] {figdata/grade-8/05-jobs-wages-unemployment/out/unemployment.csv};
\addplot table[x=year, y=jpn] {figdata/grade-8/05-jobs-wages-unemployment/out/unemployment.csv};
\legend{United States, United Kingdom, Germany, Spain, Japan}
\end{axis}
\end{tikzpicture}
\end{openingchart}

% figdata: facts (esp_max 26.1 in 2013, esp_2007 8.2; jpn_max 5.4 in 2002, jpn_1970 1.1; usa_max 9.7 in 1982;
% gbr_max 11.5 in 1985; deu_max 11.2 in 2005; usa_2020 8.0)
In Spain, the share of people who wanted a job and could not find one went from $8.2\,\%$ in 2007 to
$26.1\,\%$ in 2013: more than one in four. In Japan over the same fifty years, it never went above
$5.4\,\%$. In the United States, it rose to $9.7\,\%$ in 1982 and jumped to $8.0\,\%$ in 2020, the
year of the pandemic, from $3.7\,\%$ the year before.

What does it mean to be ``unemployed'', and who counts it? This chapter defines the words, the two
rates built on them and the survey that measures them, then looks at \omterm{def:g7:households:wage}{wages}: why they differ between
jobs, and how they have moved over a century and a half.

\begin{recall}
Hours worked and the three sectors were met in \cref{ch:g6:school-and-work}; a \omterm{def:g7:households:wage}{wage} is the pay for
work (\cref{ch:g7:households}); a \omterm{def:g8:inflation-and-price-indices:real-wage}{real wage} is a \omterm{def:g7:households:wage}{wage} divided by a \omterm{def:g8:inflation-and-price-indices:cpi}{consumer price index}, times 100
(\cref{ch:g8:inflation-and-price-indices}); output per hour measures productivity
(\cref{ch:g8:production-and-productivity}).
\end{recall}

\section{The labour market}

\begin{definition}[Labour market]\label{def:g8:jobs-wages-unemployment:labour-market}
The \emph{labour market}\index{labour market} is the set of all the exchanges in which people offer
their work and employers hire it for a \omterm{def:g7:households:wage}{wage}.
\end{definition}

Workers choose which jobs to apply for and how many hours to offer; employers choose whom to hire,
for which tasks, and at which \omterm{def:g7:households:wage}{wage}. Each day, people take new jobs, leave jobs, lose them or start
looking for one. The counts and rates below are a photograph of this \omterm{def:g7:markets-and-prices:market}{market} at one moment.

\begin{omfigure}
\includegraphics[width=0.62\linewidth]{images/book1/ai/g8-jobs-job-fair.jpg}

{\small A job fair: employers at the tables, people looking for work going from one table to the
next.}
\end{omfigure}

\section{Employed, unemployed, outside the labour force}

\begin{definition}[Employed person]\label{def:g8:jobs-wages-unemployment:employed}
An \emph{employed person}\index{employed person} is a person who did paid work in the week of the
survey, even for one hour, or who has a job and was away from it for a short time (holiday,
illness).
\end{definition}

\begin{definition}[Unemployed person]\label{def:g8:jobs-wages-unemployment:unemployed}
An \emph{unemployed person}\index{unemployed person} is a person who is not employed, is available
to start work, and has actively looked for work in the recent past.
\end{definition}

\begin{definition}[Labour force]\label{def:g8:jobs-wages-unemployment:labour-force}
The \emph{labour force}\index{labour force} is the number of \omterm{def:g8:jobs-wages-unemployment:employed}{employed persons} plus the number of
\omterm{def:g8:jobs-wages-unemployment:unemployed}{unemployed persons}.
\end{definition}

\begin{definition}[Outside the labour force]\label{def:g8:jobs-wages-unemployment:outside}
A person of working age is \emph{outside the labour force}\index{outside the labour force} when
they are neither employed nor unemployed: most students, retired people, people who care for their
family at home, and people who would like a job but have stopped looking.
\end{definition}

Being without a job is not enough to be unemployed: a retired person or a student who does not look
for work is \omterm{def:g8:jobs-wages-unemployment:outside}{outside the labour force}. A person who wants a job but has given up looking is outside
it too.

% figdata: brenland 2025 (pop 9291, under15 1490, working_age 7801, outside 3156, labour_force 4645,
% employed 4366, unemployed 279, unemp_rate 6.0, participation 59.5); under15 invented as 16 % of pop
\begin{omfigure}
\begin{tikzpicture}[x=1cm, y=1cm, font=\small]
  \node[anchor=south west, font=\small\bfseries] at (0,4.3) {Brenland, 2025: population $9\,291$};
  \draw[fill=black!8] (0,0) rectangle (2.2,3.5);
  \node[align=center] at (1.1,1.75) {under 15\\$1\,490$};
  \draw[fill=black!18] (2.2,0) rectangle (6.0,3.5);
  \node[align=center] at (4.1,1.75) {outside the\\labour force\\$3\,156$};
  \draw[fill=omDef!25] (6.0,0) rectangle (10.3,3.5);
  \node[align=center] at (8.15,1.75) {employed\\$4\,366$};
  \draw[fill=omDef!60] (10.3,0) rectangle (12,3.5);
  \node[align=center, font=\footnotesize] at (11.15,1.75) {unem-\\ployed\\$279$};
  \draw[decorate, decoration={brace, amplitude=5pt, mirror}] (2.2,-0.1) -- (12,-0.1)
    node[midway, below=6pt] {aged 15 and over: $7\,801$};
  \draw[decorate, decoration={brace, amplitude=5pt}] (6.0,3.6) -- (12,3.6)
    node[midway, above=5pt] {labour force: $4\,645$};
\end{tikzpicture}

{\small Brenland's population in 2025, in thousands, split into four groups. The number of children
under 15 is invented for this figure; the other numbers are from Brenland's accounts.}
\end{omfigure}

\section{Two rates and the survey}

\begin{definition}[Unemployment rate]\label{def:g8:jobs-wages-unemployment:unemployment-rate}
The \emph{unemployment rate}\index{unemployment rate} is the number of \omterm{def:g8:jobs-wages-unemployment:unemployed}{unemployed persons} divided by
the \omterm{def:g8:jobs-wages-unemployment:labour-force}{labour force}, in percent:
\[
\text{unemployment rate} = \frac{\text{unemployed}}{\text{employed} + \text{unemployed}} \times 100.
\]
\end{definition}

\begin{definition}[Participation rate]\label{def:g8:jobs-wages-unemployment:participation-rate}
The \emph{participation rate}\index{participation rate} is the \omterm{def:g8:jobs-wages-unemployment:labour-force}{labour force} divided by the
population of working age, in percent.
\end{definition}

In Brenland in 2025, the \omterm{def:g8:jobs-wages-unemployment:unemployment-rate}{unemployment rate} is $279 \div 4\,645 \times 100 \approx 6.0\,\%$ and the
\omterm{def:g8:jobs-wages-unemployment:participation-rate}{participation rate} is $4\,645 \div 7\,801 \times 100 \approx 59.5\,\%$. The \omterm{def:g8:jobs-wages-unemployment:unemployment-rate}{unemployment rate} is
divided by the \omterm{def:g8:jobs-wages-unemployment:labour-force}{labour force}, not by the whole population: dividing by $9\,291$ would count babies
and retired people as if they could be looking for work.

\begin{definition}[Labour force survey]\label{def:g8:jobs-wages-unemployment:survey}
A \emph{labour force survey}\index{labour force survey} is a survey that asks a sample of
\omterm{def:g7:households:household}{households}, every month or every quarter, what each of their members did in a given week, in order
to count the employed, the unemployed and the people \omterm{def:g8:jobs-wages-unemployment:outside}{outside the labour force}.
\end{definition}

\begin{measure}[The labour force survey]
Nobody is counted as unemployed because they say so. The survey asks a list of questions: did you
do any paid work last week? If not, could you have started a job? Did you do anything to find one?
In the United States, a person counts as unemployed if they were not employed in the survey week,
were available for work, and made at least one specific, active effort to find a job during the
4 weeks ending with that week; the figures cover people aged 16 and over. Other countries use the
same three tests (no work, available, searching), with small differences in ages and periods. The
international rule adds one case: a person without work who could start now and has already arranged
to start a job soon, generally within three months, counts as unemployed even without searching. The
answers of the sample are then scaled up to the whole population.
\end{measure}

\section{Why wages differ}

In the same country and the same year, some jobs pay several times more than others. Four reasons
come back again and again.

\begin{description}
  \item[What the work produces.] An employer cannot keep paying a worker more than the value of what
        that worker adds, and competing employers bid for workers who add more.
  \item[Skills and schooling.] Skills that take years to learn are scarce, and the people who have
        them add more to output.
  \item[The job itself.] Night work, danger, dirt or long trips often come with extra pay, needed to
        find people willing to do them.
  \item[Place.] \omterm{def:g7:households:wage}{Wages} in cities and in richer regions are higher, and so are \omterm{def:g7:money:price}{prices} and rents.
\end{description}

% figdata: schooling (up to 9: 163, 20430; 10 to 12: 532, 24840; 13 to 15: 559, 30770; 16 or more: 185, 40430)
Brenland's \omterm{def:g7:households:household}{household} survey of 2025 shows the second reason. Among \omterm{def:g7:households:household}{households} whose head worked
during the year, the median gross \omterm{def:g7:households:income}{income} was $20\,430$ dalens when the head had studied for 9 years
or less, $24\,840$ for 10 to 12 years, $30\,770$ for 13 to 15 years and $40\,430$ for 16 years or
more. These are \omterm{def:g7:households:household}{household} \omterm{def:g7:households:income}{incomes}, and schooling is not the only difference between these groups:
the table shows that more schooling and higher \omterm{def:g7:households:income}{income} go together, not by how much schooling alone
raises pay.

\section{Wages and productivity over a century}

\begin{model}{Wages follow productivity}
\assumes Employers compete to hire workers; none can keep paying a worker more than what that worker
adds, and a worker can leave for an employer who pays closer to it.
\predicts Over decades, \omterm{def:g8:inflation-and-price-indices:real-wage}{real wages} rise roughly as fast as \omterm{def:g7:specialisation:output-per-worker}{output per worker}.
\tested In the United States, the United Kingdom and Sweden from 1870 to 2020, \omterm{def:g8:inflation-and-price-indices:real-wage}{real wages} and real
\omterm{def:g8:gdp:gdp-per-person}{GDP per person} (a stand-in for \omterm{def:g7:specialisation:output-per-worker}{output per worker}) both rose many times over, in the same decades
(chart below).
\breaks In the short run, \omterm{def:g7:households:wage}{wages} can lag behind or run ahead of output; over long periods the two can
drift apart by a large factor, when the share of \omterm{def:g7:households:income}{income} paid as \omterm{def:g7:households:wage}{wages} changes, or when \omterm{def:g7:money:price}{prices} of
what workers buy rise faster than \omterm{def:g7:money:price}{prices} of what they produce. How large such gaps are is studied in
a later volume.
\end{model}

% figdata: real_wages 1913 = 100; 2020 values w_usa 461.8, y_usa 634.7; w_gbr 463.6, y_gbr 480.6;
% w_swe 1421.6, y_swe 936.4
\begin{omchart}{Real wages (solid) and real GDP per person (dashed), 1870--2020, both 1913 = 100, on a
times-ten scale.}{Jord\`a--Schularick--Taylor Macrohistory Database: \omterm{def:g7:households:wage}{wages} divided by consumer \omterm{def:g7:money:price}{prices};
real \omterm{def:g8:gdp:gdp-per-person}{GDP per person}}
\begin{tikzpicture}
\begin{semilogyaxis}[omtimesten, name=usa, width=0.375\linewidth, height=0.42\linewidth,
  xmin=1870, xmax=2020, ymin=20, ymax=2000, xtick={1900,1960}, title={United States\strut},
  ylabel={Index, 1913 = 100}]
\addplot[omDef, thick] table[x=year, y=w_usa] {figdata/grade-8/05-jobs-wages-unemployment/out/real_wages.csv};
\addplot[black!60, dashed] table[x=year, y=y_usa] {figdata/grade-8/05-jobs-wages-unemployment/out/real_wages.csv};
\end{semilogyaxis}
\begin{semilogyaxis}[omtimesten, name=gbr, at={(usa.east)}, anchor=west, xshift=0.3cm,
  width=0.375\linewidth, height=0.42\linewidth, xmin=1870, xmax=2020, ymin=20, ymax=2000,
  xtick={1900,1960}, title={United Kingdom\strut}, yticklabels={}]
\addplot[omDef, thick] table[x=year, y=w_gbr] {figdata/grade-8/05-jobs-wages-unemployment/out/real_wages.csv};
\addplot[black!60, dashed] table[x=year, y=y_gbr] {figdata/grade-8/05-jobs-wages-unemployment/out/real_wages.csv};
\end{semilogyaxis}
\begin{semilogyaxis}[omtimesten, at={(gbr.east)}, anchor=west, xshift=0.3cm,
  width=0.375\linewidth, height=0.42\linewidth, xmin=1870, xmax=2020, ymin=20, ymax=2000,
  xtick={1900,1960}, title={Sweden\strut}, yticklabels={}]
\addplot[omDef, thick] table[x=year, y=w_swe] {figdata/grade-8/05-jobs-wages-unemployment/out/real_wages.csv};
\addplot[black!60, dashed] table[x=year, y=y_swe] {figdata/grade-8/05-jobs-wages-unemployment/out/real_wages.csv};
\end{semilogyaxis}
\end{tikzpicture}
\end{omchart}

From 1913 to 2020, \omterm{def:g8:inflation-and-price-indices:real-wage}{real wages} were multiplied by $4.6$ in the United States and in the United
Kingdom, and by $14.2$ in Sweden; real \omterm{def:g8:gdp:gdp-per-person}{GDP per person} by $6.3$, $4.8$ and $9.4$. In all three
countries, a worker of 2020 could buy several times what a worker of 1913 could buy with a day's
pay. The two lines do not end at the same height: in the United States \omterm{def:g8:inflation-and-price-indices:real-wage}{real wages} rose less than
\omterm{def:g8:gdp:gdp-per-person}{GDP per person}, in Sweden more. The model gets the direction and the order of size right, not the
exact factor.

\section{Exercises}

\begin{exercise}[$\star$]\label{exo:g8:jobs-wages-unemployment:1}
Use the opening chart. What was Spain's highest \omterm{def:g8:jobs-wages-unemployment:unemployment-rate}{unemployment rate}, and in which year?
\end{exercise}

\begin{exercise}[$\star$]\label{exo:g8:jobs-wages-unemployment:2}
A town has $18\,800$ employed and $1\,200$ \omterm{def:g8:jobs-wages-unemployment:unemployed}{unemployed persons}. Compute its \omterm{def:g8:jobs-wages-unemployment:labour-force}{labour force} and its
\omterm{def:g8:jobs-wages-unemployment:unemployment-rate}{unemployment rate}.
\end{exercise}

\begin{exercise}[$\star$]\label{exo:g8:jobs-wages-unemployment:3}
A country has $40\,000$ people of working age, of whom $24\,000$ are in the \omterm{def:g8:jobs-wages-unemployment:labour-force}{labour force}. Compute
its \omterm{def:g8:jobs-wages-unemployment:participation-rate}{participation rate}.
\end{exercise}

\begin{exercise}[$\star$]\label{exo:g8:jobs-wages-unemployment:4}
Use the opening chart. In the 1990s, Germany's \omterm{def:g8:jobs-wages-unemployment:unemployment-rate}{unemployment rate} rose above the United Kingdom's.
About when did this happen, and about when did it fall back below it?
\end{exercise}

\begin{exercise}[$\star\star$]\label{exo:g8:jobs-wages-unemployment:5}
In 2015, Brenland's \omterm{def:g8:jobs-wages-unemployment:labour-force}{labour force} was $4\,376$ thousand and $4\,122$ thousand people were employed.
Compute the number of \omterm{def:g8:jobs-wages-unemployment:unemployed}{unemployed persons} and the \omterm{def:g8:jobs-wages-unemployment:unemployment-rate}{unemployment rate}.
\end{exercise}

\begin{exercise}[$\star\star$]\label{exo:g8:jobs-wages-unemployment:6}
In the town of exercise 2, 100 \omterm{def:g8:jobs-wages-unemployment:unemployed}{unemployed persons} stop looking for work and nothing else changes.
Compute the new \omterm{def:g8:jobs-wages-unemployment:unemployment-rate}{unemployment rate}. What happens to the \omterm{def:g8:jobs-wages-unemployment:participation-rate}{participation rate}?
\end{exercise}

\begin{exercise}[$\star\star$]\label{exo:g8:jobs-wages-unemployment:7}
Use the text. By how much were \omterm{def:g8:inflation-and-price-indices:real-wage}{real wages} multiplied from 1913 to 2020 in Sweden? And real \omterm{def:g8:gdp:gdp-per-person}{GDP per
person}?
\end{exercise}

\begin{exercise}[$\star\star$]\label{exo:g8:jobs-wages-unemployment:8}
\omterm{def:g7:specialisation:output-per-worker}{Output per worker} rises by $2\,\%$ a year for 35 years. If the model ``\omterm{def:g7:households:wage}{Wages} follow productivity''
holds, what happens to \omterm{def:g8:inflation-and-price-indices:real-wage}{real wages} over these years? (Use the rule of 70.)
\end{exercise}

\begin{exercise}[$\star\star$]\label{exo:g8:jobs-wages-unemployment:9}
Give an observation from this chapter's chart of \omterm{def:g7:households:wage}{wages} that does not fit the model ``\omterm{def:g7:households:wage}{Wages} follow
productivity'' exactly.
\end{exercise}

\begin{exercise}[$\star\star\star$]\label{exo:g8:jobs-wages-unemployment:10}
A hospital employs 300 nurses paid $3\,000$ dalens a month and 100 doctors paid $7\,000$ dalens a
month. Compute the average monthly \omterm{def:g7:households:wage}{wage} in the hospital.
\end{exercise}

\begin{exercise}[$\star\star\star$]\label{exo:g8:jobs-wages-unemployment:11}
A country has $50\,000$ people of working age, $28\,500$ employed and $1\,500$ unemployed. (a) Compute
the \omterm{def:g8:jobs-wages-unemployment:unemployment-rate}{unemployment rate} and the \omterm{def:g8:jobs-wages-unemployment:participation-rate}{participation rate}. (b) 500 of the unemployed stop looking for work;
employment does not change. Compute both rates again. (c) Did anyone find a job?
\end{exercise}

\begin{exercise}[$\star\star\star$]\label{exo:g8:jobs-wages-unemployment:12}
Give two things about work that the \omterm{def:g8:jobs-wages-unemployment:unemployment-rate}{unemployment rate} cannot show.
\end{exercise}

\section{Problem: The Brenland Labour Survey}

\begin{problem}[{Weekend problem --- classifying survey answers, computing the two rates, and
Brenland's unemployment from 2000 to 2025}]
\label{pb:g8:jobs-wages-unemployment:1}
Everything in this problem is invented. A small test of Brenland's \omterm{def:g8:jobs-wages-unemployment:survey}{labour force survey} asks 20
people about the week of 3 to 9 March.
\begin{center}
\small
\begin{tabular}{rrl@{\hspace{1.5em}}rrl}
\toprule
No. & Age & Answer & No. & Age & Answer \\
\midrule
1 & 34 & worked 40 hours & 11 & 41 & worked 38 hours \\
2 & 16 & at school; not looking & 12 & 60 & no job; looking; can start now \\
3 & 71 & retired & 13 & 33 & home with children; not looking \\
4 & 45 & no job; looking; can start now & 14 & 27 & worked 45 hours \\
5 & 29 & worked 10 hours & 15 & 8 & at school \\
6 & 12 & at school & 16 & 48 & has a job; on holiday \\
7 & 52 & no job; stopped looking & 17 & 36 & no job; looking; cannot start before June \\
8 & 38 & worked 35 hours & 18 & 55 & worked 40 hours \\
9 & 23 & no job; looking; can start now & 19 & 25 & no job; looking; can start now \\
10 & 19 & student; worked 8 hours & 20 & 66 & worked 15 hours \\
\bottomrule
\end{tabular}
\end{center}

\medskip
\noindent\textbf{Part I --- Classifying.}
\begin{enumerate}
  \item Put each person in one group: under 15, employed, unemployed, or \omterm{def:g8:jobs-wages-unemployment:outside}{outside the labour force}.
  \item How many people are in each group?
  \item Person 7 would like a job. Why is she not unemployed?
  \item Why is person 17 not unemployed, although he is looking for work?
\end{enumerate}

\medskip
\noindent\textbf{Part II --- The rates.}
\begin{enumerate}[resume]
  \item What is the \omterm{def:g8:jobs-wages-unemployment:labour-force}{labour force} of the sample?
  \item Compute the \omterm{def:g8:jobs-wages-unemployment:unemployment-rate}{unemployment rate} of the sample.
  \item Compute the \omterm{def:g8:jobs-wages-unemployment:participation-rate}{participation rate} of the sample (working age is 15 and over).
  \item Person 7 starts looking again and could start at once. Compute both rates again.
\end{enumerate}

\medskip
\noindent\textbf{Part III --- Brenland, 2000--2025.} The table gives Brenland's \omterm{def:g8:jobs-wages-unemployment:labour-force}{labour force} and
employment, in thousands.
\begin{center}
\begin{tabular}{lrrrrrrr}
\toprule
Year & 2000 & 2008 & 2009 & 2012 & 2020 & 2021 & 2025 \\
\midrule
\omterm{def:g8:jobs-wages-unemployment:labour-force}{Labour force} & 4\,000 & 4\,196 & 4\,221 & 4\,298 & 4\,508 & 4\,535 & 4\,645 \\
Employed & 3\,760 & 3\,906 & 3\,883 & 4\,019 & 4\,111 & 4\,317 & 4\,366 \\
\bottomrule
\end{tabular}
\end{center}
\begin{enumerate}[resume]
  \item Compute the \omterm{def:g8:jobs-wages-unemployment:unemployment-rate}{unemployment rate} in 2000 and in 2025.
  \item Compute it in 2008 and in 2009. How many more people were unemployed in 2009?
  \item In 2025, $7\,801$ thousand people in Brenland were aged 15 and over. Compute the
        \omterm{def:g8:jobs-wages-unemployment:participation-rate}{participation rate}.
  \item Among the years of the table, which had the highest \omterm{def:g8:jobs-wages-unemployment:unemployment-rate}{unemployment rate}, and what was it?
\end{enumerate}
\end{problem}
